\documentclass[a4paper]{article}

% --- Основные пакеты ---
\usepackage{graphicx}     % картинки
\usepackage{geometry}     % поля страницы
\usepackage{hyperref}     % гиперссылки
\usepackage{amsmath}      % математика
\usepackage{amssymb}      % математические символы
\usepackage{float}        % [H] позиционирование фигур
\usepackage[T2A]{fontenc} % кириллица
\usepackage[utf8]{inputenc} % кодировка UTF-8
\usepackage[russian]{babel} % многоязычность
\usepackage{caption}      % настройка подписей

% --- Настройка страницы ---
\geometry{top=1.5cm, bottom=1.5cm, left=3cm, right=1.5cm}

% --- Настройка гиперссылок ---
\hypersetup{
    colorlinks,
    linkcolor=black,
    urlcolor=blue,
    citecolor=black
}

% --- Настройка подписей ---
\renewcommand{\figurename}{Рис.}
\captionsetup{labelsep=period}

\begin{document}

% --- Титульная страница ---
\centering
Министерство науки и высшего образования Российской Федерации\\
Федеральное государственное автономное образовательное учреждение\\
высшего образования\\
«Южно-Уральский государственный университет\\
(национальный исследовательский университет)»\\
Институт естественных и точных наук

\setlength{\parskip}{7.5cm}

Отчет\\
Лабораторная работа №6.\\
Методы численного интегрирования.

\setlength{\parskip}{6.5cm}

\raggedleft
\begin{minipage}{0.345\textwidth}
\raggedright
Выполнил: студент группы ЕТ-210\\
Тельнов С.М.\\
Проверил: преподаватель кафедры ВМ\\
Стариков Я.Е.
\end{minipage}

\setlength{\parskip}{6cm}

\centering
Челябинск — 2026

\setlength{\parskip}{0pt plus 1pt}
\newpage

\setlength{\parskip}{0pt plus 1pt}

% --- Содержание ---
\centering
\textbf{Краткое содержание}
\begin{itemize}
    \setlength{\parskip}{0.4cm}
    \item \hyperref[sec:theory]{\text{Теоретическая часть} \dotfill \pageref{sec:theory}}
    \item \hyperref[sec:practice]{\text{Практическая часть} \dotfill \pageref{sec:practice}}
    \item \hyperref[sec:conclusion]{\text{Вывод} \dotfill \pageref{sec:conclusion}}
\end{itemize}

\setlength{\parindent}{1.5em}

\raggedright

% --- Теоретическая часть ---
\section*{\centering\fontsize{14pt}{0pt}\selectfont\textbf{Теоретическая часть}\label{sec:theory}}

\section*{\fontsize{11pt}{0pt}\selectfont\textbf{Методы численного интегрирования}}
Задача численного интегрирования заключается в нахождении приближенного значения определенного интеграла $I = \int_a^b f(x) dx$.
Отрезок интегрирования $[a, b]$ разбивается на $N$ равных частичных отрезков с шагом $h = \frac{b-a}{N}$, где узлы сетки задаются как $x_i = a + i \cdot h$, $i = 0, 1, \dots, N$.

\subsubsection*{(a) Метод прямоугольников}
Метод основан на замене подынтегральной функции кусочно-постоянной функцией.
\begin{itemize}
    \item \textbf{Левые прямоугольники:} значение функции берется в левой границе каждого отрезка: 
    $I \approx h \sum_{i=0}^{N-1} f(x_i)$.
\item \textbf{Правые прямоугольники:} значение функции берется в правой границе: 
    $I \approx h \sum_{i=1}^{N} f(x_i)$.
\item \textbf{Центральные прямоугольники:} значение функции берется в середине отрезка $x_{i+1/2} = \frac{x_i + x_{i+1}}{2}$: 
    $I \approx h \sum_{i=0}^{N-1} f(x_{i+1/2})$.
Этот метод обладает более высоким порядком точности по сравнению с левыми и правыми прямоугольниками.
\end{itemize}

\subsubsection*{(b) Метод трапеций}
На каждом интервале $[x_i, x_{i+1}]$ функция аппроксимируется линейным полиномом (хордой).
Площадь под графиком заменяется суммой площадей трапеций:
$$I \approx \frac{h}{2} \sum_{i=0}^{N-1} (f(x_i) + f(x_{i+1})) = h \left( \frac{f(a)+f(b)}{2} + \sum_{i=1}^{N-1} f(x_i) \right)$$

\subsubsection*{(c) Метод Симпсона (метод парабол)}
На каждой паре соседних интервалов функция аппроксимируется полиномом второй степени (параболой).
Требует четного количества отрезков разбиения ($N = 2k$):
$$I \approx \frac{h}{3} \left[ f(x_0) + 4\sum_{i=1, 3, \dots}^{N-1} f(x_i) + 2\sum_{j=2, 4, \dots}^{N-2} f(x_j) + f(x_N) \right]$$

\subsubsection*{(d) Квадратуры Гаусса}
Метод Гаусса обеспечивает наивысшую алгебраическую точность $2m-1$ для $m$ узлов.
Интеграл приводится к стандартному интервалу $[-1, 1]$ заменой $x = \frac{b-a}{2}t + \frac{a+b}{2}$:
$$I = \frac{b-a}{2} \int_{-1}^1 f\left(\frac{b-a}{2}t + \frac{a+b}{2}\right) dt \approx \frac{b-a}{2} \sum_{k=1}^m w_k f(\tilde{x}_k)$$
Узлы $t_k$ являются корнями полинома Лежандра степени $m=\{3,5,7\}$, а веса $w_k$ определяются интегрированием базисных полиномов Лагранжа.

\subsubsection*{(e) Метод Монте-Карло}
Стохастический метод вычисления интеграла. Для многомерного случая $I = \int_D f(x) dx$, где $D$ — область объема $V$, интеграл оценивается как среднее значение функции в случайных точках $X_k$, равномерно распределенных в $D$:
$$I \approx \frac{V}{N} \sum_{k=1}^N f(X_k)$$
Ошибка метода убывает пропорционально $1/\sqrt{N}$, что делает его эффективным для высокой размерности.

---

\section*{\fontsize{11pt}{0pt}\selectfont\textbf{Сведение многомерных интегралов к повторным и аналитическое решение}}

\subsubsection*{(a) Объем 3-мерного и 5-мерного единичного шара}
Сведение к повторному интегралу и аналитическое решение для 3-мерного шара (с переходом к сферическим координатам):
$$V_3 = \int_{-1}^1 dx \int_{-\sqrt{1-x^2}}^{\sqrt{1-x^2}} dy \int_{-\sqrt{1-x^2-y^2}}^{\sqrt{1-x^2-y^2}} dz = \int_0^{2\pi} d\varphi \int_0^\pi \sin\theta d\theta \int_0^1 r^2 dr = 2\pi \cdot 2 \cdot \frac{1}{3} = \frac{4}{3}\pi \approx 4.1888$$

Для 5-мерного шара применим рекуррентную формулу объема $n$-мерной сферы $V_n = \frac{2\pi}{n}V_{n-2}$:
$$V_5 = \int_{-1}^1 dx_1 \dots \int_{-\sqrt{1-\sum_{i=1}^4 x_i^2}}^{\sqrt{1-\sum_{i=1}^4 x_i^2}} dx_5 = \frac{2\pi}{5} V_3 = \frac{2\pi}{5} \cdot \frac{4\pi}{3} = \frac{8\pi^2}{15} \approx 5.2638$$

\subsubsection*{(b) Вычисление двумерного интеграла}
Задан интеграл $\iint_{D}(x^{2}+y^{2})dxdy$, где область $D$ ограничена параболами $y = x^2 - 1$ и $y = 1 - x^2$.
Найдем точки пересечения: $x^2 - 1 = 1 - x^2 \implies 2x^2 = 2 \implies x = \pm 1$.

Сведение к повторному интегралу:
$$I = \int_{-1}^1 dx \int_{x^2-1}^{1-x^2} (x^2+y^2) dy$$
Вычислим внутренний интеграл:
$$\int_{x^2-1}^{1-x^2} (x^2+y^2) dy = \left[ x^2y + \frac{y^3}{3} \right]_{x^2-1}^{1-x^2} = 2 \left( x^2(1-x^2) + \frac{(1-x^2)^3}{3} \right) = \frac{2}{3} (1 - x^6)$$
Вычислим внешний интеграл:
$$I = \frac{2}{3} \int_{-1}^1 (1 - x^6) dx = \frac{4}{3} \int_0^1 (1 - x^6) dx = \frac{4}{3} \left[ x - \frac{x^7}{7} \right]_0^1 = \frac{4}{3} \cdot \frac{6}{7} = \frac{8}{7} \approx 1.1429$$

\subsubsection*{(c) Вычисление объема конуса}
Задан интеграл $\iiint_{D} dxdydz$, где $D: x^2+y^2 \le z^2 \tan(\theta)$, $\theta = \pi/6$, $0 \le z \le 2$.

Сведение к повторному интегралу в цилиндрических координатах ($x = r\cos\varphi$, $y = r\sin\varphi$, $z = z$):
$$V = \int_0^{2\pi} d\varphi \int_0^2 dz \int_0^{z\sqrt{\tan(\pi/6)}} r dr$$
Выполним аналитическое интегрирование:
$$V = 2\pi \int_0^2 \left[ \frac{r^2}{2} \right]_0^{z\sqrt{\tan(\pi/6)}} dz = \pi \tan(\pi/6) \int_0^2 z^2 dz = \frac{\pi}{\sqrt{3}} \left[ \frac{z^3}{3} \right]_0^2 = \frac{8\pi}{3\sqrt{3}} \approx 4.8368$$

\subsubsection*{(d) Вычисление трехмерного интеграла с тригонометрией}
Задан интеграл $I = \iiint_{D}\cos(x)\cos(y)\cos(z)dxdydz$ по той же области $D$.
Сведение к повторному интегралу с использованием цилиндрических координат:
$$I = \int_0^2 \cos(z) dz \iint_{x^2+y^2 \le z^2/\sqrt{3}} \cos(x)\cos(y) dx dy$$
В полярных координатах внутренний двойной интеграл по кругу радиуса $R = z/\sqrt[4]{3}$ математически выражается через функцию Бесселя первого рода первого порядка $J_1$, что дает:
$$\iint_{x^2+y^2 \le R^2} \cos(x)\cos(y) dx dy = \sqrt{2}\pi R J_1(\sqrt{2}R)$$
Подставляя радиус $R$, сводим задачу к вычислению финального однократного определенного интеграла:
$$I = \int_0^2 \cos(z) \left( \sqrt{2}\pi \frac{z}{\sqrt[4]{3}} J_1\left(\sqrt{2}\frac{z}{\sqrt[4]{3}}\right) \right) dz \approx 0.4037$$

\section*{\centering\fontsize{14pt}{0pt}\selectfont\textbf{Практическая часть}\label{sec:practice}}

\subsection*{Задание 1: Использование методов для расчётов}

На первом этапе было вычислено значение определённого интеграла различными численными методами. Результаты представлены в таблице:

\begin{table}[H]
\centering
\begin{tabular}{lr}
\hline
Метод & Результат \\
\hline
Правый прямоугольник & 37.0757 \\
Левый прямоугольник & 37.1716 \\
Центральный прямоугольник & 37.1237 \\
Трапеция & 37.1237 \\
Симпсон & 37.1237 \\
Гаусс & 35.9095 \\
Монте-Карло & 38.9662 \\
Метод из scipy.integrate & 37.1237 \\
Аналитический & 37.1237 \\
\hline
\end{tabular}
\caption{Результаты численного интегрирования различными методами (задание 1)}
\end{table}

Из таблицы видно, что наиболее точные результаты дают методы Симпсона, трапеций и центральных прямоугольников, которые совпадают с аналитическим решением. Метод Гаусса дал значительную погрешность, вероятно из-за неправильно выбранных узлов или веса.

\subsection*{Задание 2: Сравнение методов с разными N}

Сравнено влияние числа разбиений $N$ на точность вычислений для различных методов интегрирования. Аналитическое решение интеграла равно 37.1237.

\begin{table}[H]
\centering
\small
\begin{tabular}{lrr}
\hline
Метод & N & Результат \\
\hline
Правый прямоугольник & 100 & 36.6414 \\
& 1000 & 37.0757 \\
Левый прямоугольник & 100 & 37.6003 \\
& 1000 & 37.1716 \\
Центральный прямоугольник & 100 & 37.1251 \\
& 1000 & 37.1237 \\
Трапеция & 100 & 37.1209 \\
& 1000 & 37.1237 \\
Симпсон & 100 & 37.1237 \\
& 1000 & 37.1237 \\
\hline
\end{tabular}
\caption{Значения интеграла различными методами при различных N}
\label{tab:n_comparison}
\end{table}

Таблица абсолютных ошибок при сравнении с аналитическим решением:

\begin{table}[H]
\centering
\small
\begin{tabular}{lrr}
\hline
N & Метод & Абсолютная ошибка \\
\hline
100 & Правый прямоугольник & $4.8234 \times 10^{-1}$ \\
1000 & Правый прямоугольник & $4.7978 \times 10^{-2}$ \\
100 & Левый прямоугольник & $4.7665 \times 10^{-1}$ \\
1000 & Левый прямоугольник & $4.7921 \times 10^{-2}$ \\
100 & Центральный прямоугольник & $1.4222 \times 10^{-3}$ \\
1000 & Центральный прямоугольник & $1.4214 \times 10^{-5}$ \\
100 & Трапеция & $2.8437 \times 10^{-3}$ \\
1000 & Трапеция & $2.8428 \times 10^{-5}$ \\
100 & Симпсон & $3.5838 \times 10^{-6}$ \\
1000 & Симпсон & $3.5757 \times 10^{-10}$ \\
\hline
\end{tabular}
\caption{Абсолютные ошибки различных методов при различных N}
\label{tab:errors}
\end{table}

Для метода Монте-Карло исследовано влияние числа итераций $N$ на точность:

\begin{table}[H]
\centering
\small
\begin{tabular}{lrr}
\hline
N & Результат & Абсолютная ошибка \\
\hline
1000 & 37.2775 & $1.5385 \times 10^{-1}$ \\
10000 & 37.3200 & $1.9628 \times 10^{-1}$ \\
100000 & 37.1363 & $1.2644 \times 10^{-2}$ \\
\hline
\end{tabular}
\caption{Результаты метода Монте-Карло при различных числах итераций}
\label{tab:mc_results}
\end{table}

\subsection*{Задание 3: Вычисление многомерных интегралов}

Вычислены различные многомерные интегралы всеми доступными численными методами.

\subsubsection*{Интеграл (a): Объем единичного шара}

Для 3-мерного случая объём единичного шара вычисляется аналитически как $V = \frac{4}{3}\pi \approx 4.1888$. Результаты численных методов представлены в таблице:

\begin{table}[H]
\centering
\small
\begin{tabular}{lrr}
\hline
Метод & Результат & Абсолютная погрешность \\
\hline
Правый прямоугольник & 4.1750 & $1.3819 \times 10^{-2}$ \\
Левый прямоугольник & 4.1750 & $1.3827 \times 10^{-2}$ \\
Центральный прямоугольник & 4.1931 & $4.3054 \times 10^{-3}$ \\
Трапеция & 4.1750 & $1.3827 \times 10^{-2}$ \\
Симпсон & 4.1839 & $4.8822 \times 10^{-3}$ \\
Гаусс & 4.1941 & $5.2852 \times 10^{-3}$ \\
Монте-Карло & 4.1822 & $6.5560 \times 10^{-3}$ \\
\hline
\end{tabular}
\caption{Вычисление объёма 3-мерного единичного шара различными методами}
\label{tab:sphere_3d}
\end{table}

Для 5-мерного случая аналитическое решение равно $V = \frac{8\pi^2}{15} \approx 5.2638$:

\begin{table}[H]
\centering
\small
\begin{tabular}{lrr}
\hline
Метод & Результат & Абсолютная погрешность \\
\hline
Правый прямоугольник & 5.2384 & $2.5374 \times 10^{-2}$ \\
Левый прямоугольник & 5.2384 & $2.5374 \times 10^{-2}$ \\
Центральный прямоугольник & 5.2719 & $8.0862 \times 10^{-3}$ \\
Трапеция & 5.2384 & $2.5374 \times 10^{-2}$ \\
Симпсон & 5.2551 & $8.6620 \times 10^{-3}$ \\
Гаусс & 5.2699 & $6.0984 \times 10^{-3}$ \\
Монте-Карло & 5.2410 & $2.2846 \times 10^{-2}$ \\
\hline
\end{tabular}
\caption{Вычисление объёма 5-мерного единичного шара различными методами}
\label{tab:sphere_5d}
\end{table}

\subsubsection*{Интеграл (b): Двумерный интеграл}

Аналитическое решение равно $I = \frac{8}{7} \approx 1.1429$. Результаты вычислений в трёхмерном случае:

\begin{table}[H]
\centering
\small
\begin{tabular}{lrr}
\hline
Метод & Результат & Абсолютная погрешность \\
\hline
Правый прямоугольник & 1.1429 & $4.6391 \times 10^{-5}$ \\
Левый прямоугольник & 1.1429 & $4.6391 \times 10^{-5}$ \\
Центральный прямоугольник & 1.1429 & $4.1092 \times 10^{-5}$ \\
Трапеция & 1.1429 & $4.6391 \times 10^{-5}$ \\
Симпсон & 1.1429 & $4.2857 \times 10^{-5}$ \\
Гаусс & 1.1429 & $4.2857 \times 10^{-5}$ \\
Монте-Карло & 1.1392 & $3.7364 \times 10^{-3}$ \\
\hline
\end{tabular}
\caption{Вычисление интеграла (b) различными методами в трёхмерном случае}
\label{tab:int_b_3d}
\end{table}

\subsubsection*{Интеграл (c): Объем конуса}

Аналитическое решение равно $V = \frac{8\pi}{3\sqrt{3}} \approx 4.8368$. Результаты вычислений:

\begin{table}[H]
\centering
\small
\begin{tabular}{lrr}
\hline
Метод & Результат & Абсолютная погрешность \\
\hline
Правый прямоугольник & 4.8513 & $1.4518 \times 10^{-2}$ \\
Левый прямоугольник & 4.8223 & $1.4502 \times 10^{-2}$ \\
Центральный прямоугольник & 4.8368 & $6.5322 \times 10^{-6}$ \\
Трапеция & 4.8368 & $7.9782 \times 10^{-6}$ \\
Симпсон & 4.8368 & $1.6954 \times 10^{-6}$ \\
Гаусс & 4.8368 & $1.6954 \times 10^{-6}$ \\
Монте-Карло & 4.8276 & $9.2280 \times 10^{-3}$ \\
\hline
\end{tabular}
\caption{Вычисление объёма конуса различными методами}
\label{tab:cone}
\end{table}

\subsubsection*{Интеграл (d): Трёхмерный интеграл с тригонометрией}

Результаты вычислений методом Монте-Карло показали значительную погрешность:

\begin{table}[H]
\centering
\small
\begin{tabular}{lr}
\hline
Метод & Результат \\
\hline
Правый прямоугольник & 0.5934 \\
Левый прямоугольник & 0.5979 \\
Центральный прямоугольник & 0.5957 \\
Трапеция & 0.5956 \\
Симпсон & 0.5957 \\
Гаусс & 0.6058 \\
Монте-Карло & 0.4348 \\
\hline
\end{tabular}
\caption{Вычисление интеграла (d) различными методами}
\label{tab:int_d}
\end{table}


\section*{\centering\fontsize{14pt}{0pt}\selectfont\textbf{Вывод}\label{sec:conclusion}}

Проведённые вычисления демонстрируют эффективность различных методов численного интегрирования. Методы Симпсона, трапеций и центральных прямоугольников показали наивысшую точность для одномерного интеграла, где абсолютная ошибка при $N=1000$ достигала порядка $10^{-5}$ — $10^{-10}$. Это подтверждает теоретический порядок точности методов: трапеции имеют погрешность $O(h^2)$, центральные прямоугольники $O(h^2)$, а метод Симпсона $O(h^4)$.

Для многомерных интегралов зависимость от размерности становится более сложной. Методы Монте-Карло демонстрируют преимущество в высоких размерностях благодаря независимости ошибки от степени мультипликативного роста при увеличении числа измерений. Однако для малых размерностей (2–3 измерения) детерминированные методы дают более точные результаты.

Особое внимание следует уделить выбору узлов квадратур Гаусса, который критически влияет на точность метода. Результаты показывают, что при неправильном выборе параметров метод может давать погрешность, превышающую другие методы в разы.
\end{document}
